Mixed convection in a rotating eccentric horizontal annulus was studied using 5,460 numerical simulations spanning 35 geometries and a joint Reynolds number-Richardson number parameter space, replacing the classical Rayleigh number and rotation rate convention, justified on physical grounds for a single rotating boundary. The dataset was validated against experimental, numerical, and analytical benchmarks from the literature, with a mean absolute difference of $5.3\%$ against sixteen exactly matched experimental points. Three candidate architectures were compared as scalar Nusselt number surrogates; a gradient-boosted tree ensemble was most accurate, reaching $R^2=0.9996$ on a held-out split of the training grid. Because same-grid accuracy can overstate real-world reliability, the surrogate was tested against eighteen new simulations at operating conditions absent from the training data: this true out-of-sample test gives a mean error of $7.89\%$, and shows that the model's own uncertainty estimate, well calibrated on the training grid, becomes unreliable under genuine extrapolation, with coverage of its own 95\% interval falling to $27.8\%$. Interpretability analysis recovered physically consistent effect directions for every input, purely from the trained model. For the spatial temperature and velocity field, a per-node network outperformed both a DeepONet and a graph neural network built on the true mesh connectivity, the same ranking recovered independently on this dataset as under an earlier, differently parameterized version of this study. These findings argue for routinely validating surrogate models against genuinely new simulations, not only held-out splits of the training grid, and for treating ensemble-based uncertainty estimates cautiously outside the training distribution.
